\section{Introduction}

Land cover data products provide valuable information about the characteristics of the Earth's surface for a wide variety of applications, including agriculture, urban planning, and environmental monitoring~\citep{scarpace1980land, lovelandDevelopmentGlobalLand2000,townshendLandCover1992,wulderLandCover2018,banGlobalLandCover2015}.
Research into automated methods for land cover classification from remotely sensed data has been ongoing for several decades, with a majority of methods relying on supervised learning techniques to classify land cover using multispectral aerial or satellite imagery~\citep{phiriDevelopmentsLandsatLand2017, maxwellImplementationMachinelearningClassification2018}.
Historically, pixel-wise classification methods have been the predominant approach for land cover classification, where each pixel in the image is classified based on the pixel's raw spectral values, hand-crafted features, and ancillary data from other geospatial datasets~\citep{mciverEstimatingPixelscaleLand2001, khatamiMetaanalysisRemoteSensing2016}.
This pixel-wise paradigm is effective for classifying coarse and moderate spatial resolution imagery (e.g., MODIS, Landsat, Sentinel-2), but is less effective for classifying high spatial resolution imagery due to decreased inter-class spectral separability and increased intra-class spectral variability~\citep{strahlerNatureModelsRemote1986,woodcockFactorScaleRemote1987,hansenReviewLargeArea2012, grekousisOverview21Global2015, markhamLandCoverClassification1981, luSurveyImageClassification2007}.
This is problematic as very high spatial resolution (VHSR, \(<\) \SI{5}{\meter}) land cover data offer more detailed information about the Earth's surface with greater spatial accuracy, which is valuable for applications that require fine-grained information about land cover types and their spatial distribution~\citep{liImpactsSpatialResolutions2022,fisherImpactSatelliteImagery2018}.
Object-based Image Analysis (OBIA) techniques have been proposed as a potential workflow to address these issues, but their performance is highly dependent on the quality of objects or superpixels extracted from the imagery during the segmentation process, which must be carefully tuned to the specific characteristics of the imagery and the land cover classes of interest~\citep{blaschkeGeographicObjectBasedImage2014a, johnsonUnsupervisedImageSegmentation2011, maReviewSupervisedObjectbased2017}.
Nevertheless, OBIA approaches demonstrate the value of leveraging spatial relationships in the imagery to provide more information to the classifier, directly addressing the inter-class spectral separability issues that hamper pixel-wise classification methods~\citep{blaschkeObjectOrientedImageProcessing2000a, myintPerpixelVsObjectbased2011, tehranyComparativeAssessmentObject2014}.
The application of deep computer vision models to land cover classification has been of significant interest in recent years, as these models excel at extracting spatial and spectral features from imagery in a data-driven manner, making them well-suited for high spatial resolution land cover classification tasks~\citep{zhuDeepLearningRemote2017a, maDeepLearningRemote2019, valiDeepLearningLand2020}.


% Semantic segmentation
Convolutional neural networks (CNNs) are a class of deep learning models that have been widely used for land cover classification tasks due to their ability to automatically extract hierarchical spatial structures from the input imagery that are useful for classifying land cover type~\citep{liSurveyConvolutionalNeural2022a, songPatchBasedLightConvolutional2019, panLandcoverClassificationMultispectral2020,kussulDeepLearningClassification2017}.
% Semantic segmentation models are a class of CNNs that consist of two main components: an encoder network that extracts spatial and spectral features from the input imagery, and a decoder network that generates a pixel-wise classification map from the features extracted by the encoder.
CNN-based semantic segmentation models learn a mapping from the input image to a pixel-wise classification probability map, which can be used to generate a land cover map by assigning each pixel to the class with the highest probability~\citep{longFullyConvolutionalNetworks2015}.
Semantic segmentation models have been widely used for land cover classification tasks, as they can facilitate the generation of high spatial resolution land cover maps that leverage both spectral and spatial information in the input imagery under a unified framework~\citep{yuanReviewDeepLearning2021}.
However, these models require large amounts of labeled training data to learn the complex spatial and spectral patterns present in the imagery, which can be difficult to obtain for high spatial resolution imagery due to the high cost of collecting and labeling the data~\citep{huTransferringDeepConvolutional2015}.

This need for large amounts of ground truth data for accurate training of deep learning models has led to the extensive research into methods for reducing the amount of data necessary to train these models.
Most of these methods focus on leveraging transfer learning, where the parameters of a model pre-trained using a large dataset are transferred to a model that is subsequently fine-tuned on a smaller dataset of interest.
Transfer learning is effective for deep computer vision models, as the features learned by a model on a large dataset are often generalizable to other datasets thanks to the hierarchical nature of the features learned by the model~\citep{huTransferringDeepConvolutional2015, kolesnikovBigTransferBiT2020}.
However, the effectiveness of transfer learning is highly dependent on the similarity between the source and target datasets, as well as the amount of labeled data available for the target dataset.
For example, a common approach to transfer learning for land cover classification tasks is to pre-train a model on a large general-purpose image classification dataset (e.g., ImageNet~\citep{russakovskyImageNetLargeScale2015}) and then fine-tune the model on a smaller dataset of interest~\citep{piresdelimaConvolutionalNeuralNetwork2020}.
This is generally effective, but the domain shift between the source and target datasets can hinder the performance of the model, and the extent of this domain shift between general purpose image classification datasets and high spatial resolution remote sensing datasets is not well understood~\citep{maTransferLearningEnvironmental2024}.
Self-supervised pre-training methods have been proposed as an alternative to traditional fully supervised pre-training methods, as they can leverage the spatial and spectral information present in the input imagery to learn useful features through the use of pretext tasks that do not require labeled data~\citep{wuUnsupervisedFeatureLearning2018,tianContrastiveMultiviewCoding2020, chenSimpleFrameworkContrastive2020}.
% These methods have quickly gained popularity in the 
These approaches are rapidly gaining traction as a means to pre-train large models using medium-to-high satellite imagery that can serve as strong feature extraction backbones for a variety of downstream tasks in remote sensing analysis without the need for ground truth data.
Self-supervised pre-training methods are particularly attractive for remote sensing applications, as large amounts of labeled imagery datasets are often unavailable, but source imagery is plentiful thanks to the proliferation of Earth observation satellites and aerial imagery platforms~\citep{xiaoFoundationModelsRemote2025}.
However, few studies exist that rigorously investigate the applicability of these methods to operational land cover mapping at VHSR under extreme low-data scenarios, though similar works into their data efficiency have yielded promising results in moderate resolution land cover classification tasks~\citep{ayushGeographyAwareSelfSupervisedLearning2021, manasSeasonalContrastUnsupervised2021a}.

% However, no works on applying foundation models to vision 
% However, foundation models are typically oriented for cross-task, multi-sensor, multi-modal, or multi-temporal analysis scenarios, and while they often claim to be label-efficient, few works rigorously investigate their performance under extreme low-data scenarios for specific tasks such as land cover classification at VHSR~\citep{zhangSelfsupervisedLearningRemote2024, zhuSelfSupervisedLearning2023}.
% While effective, foundation models are geared for performance in multi-sensor, mutli-modal, or cross-task scenarios, and while claims of 
% As such, to the authors' knowledge, little work has been done on how these self-supervised pre-training approaches can impact operational land cover mapping under extreme low-data scenarios.
% These methods are effective, but most studies on the effectiveness of self-supervised learning are focused on single-stage pre-training methods where the only method used for pre-training is self-supervised learning, robbing the model of the benefits of supervised pre-training~\citep{dippelFinegrainedVisualRepresentations2022,manasSeasonalContrastUnsupervised2021a, ayushGeographyAwareSelfSupervisedLearning2021, scheibenreifSelfsupervisedVisionTransformers2022}.

Beyond the choice of pre-training routine, another poorly understood component of developing semantic segmentation models for land cover classification is the choice of sampling routine used to select training data.
Ideally, a simple random sample of sufficient size from the region of interest should produce a quality training dataset.
However, the sample size must be large enough to satisfy the conditions that the training data samples are representative of the variation in land cover across the region while also having enough samples in the minority classes such that a classifier can capture said classes accurately.
This makes simple random sampling problematic for large areas with diverse land cover types or strongly imbalanced land cover distributions.
Much research on stratification strategies for building pixel-based datasets for land cover classification models has been conducted to address the issues associated with simple random sampling~\citep{ zhenImpactTrainingValidation2013,canoImprovedForestcoverMapping2017, stehmanImpactSampleSize2012}.
These stratification procedures involve splitting the population into several groups based on some intrinsic property (e.g., geospatial location) or information from an external data source (e.g., prior land cover information) and sampling from these groups proportionally~\citep{costaSpatiallyStratifiedMultiStage2022, wangMeasureSpatialStratified2016,freemanEvaluatingEffectivenessDownsampling2012}.
However, semantic segmentation models require small patches consisting of dense labels, that is, every pixel in the patch is assigned to a class, rather than individual pixels.
While this does not negate the implications of stratification, it does add extra considerations to the process, as one patch can have multiple classes.
For example, when stratifying training data by land cover class according to some prior knowledge of the land cover distribution, it can be very difficult to select patches in such a way that fit the desired land cover distribution.
To date, no works exist that explicitly investigate stratification strategies through the lens of patch-based training samples for semantic segmentation models for land cover classification, despite the importance of quality training data selection for said models.
% Further, the composition of classes in sampled patches should also be representative of spatial relationships between land cover types in the region of interest. \note{SEE} \url{https://chatgpt.com/c/6894f587-a918-832a-8112-73ab5926bb09}, \url{https://www.mdpi.com/2072-4292/15/6/1607#:~:text=Figure%205C%2C%20we%20show%20that,the%20reported%20confusion%20matrix%20accuracies}
% This issue is exacerbated under low-data scenarios, where it can be difficult to capture the full diversity of both land cover types and their spatial relationships with few samples.

In this work, we address the challenge of operational land cover classification at VHSR under label-scarce conditions by introducing and validating two complementary strategies: a novel stratification approach for selecting training samples and the use of self-supervised learning for pre-training deep convolutional encoder networks.
We demonstrate that employing a semantic stratification strategy that groups candidate patches according to their semantic content (i.e., land cover composition) using a coarse-resolution land cover dataset to guide the stratification selection to sample patches that are representative of the diversity of land cover compositions found in the region of interest.
Additionally, we employ Bootstrap Your Own Latent (BYOL) as a self-supervised pretext task during pre-training, enabling a ResNet-101 convolutional encoder network to learn to extract deep features from aerial imagery without the need for full supervision using labeled data~\citep{grillBootstrapYourOwn2020}.
This approach is used to tackle land cover classification at \SI{1}{\meter} resolution over the state of Mississippi using only 1,000 patches selected using the semantic stratification procedure.
We demonstrate that these approaches can be used to train deep semantic segmentation models that are accurate, generalizable, and capable of yielding acceptable performance even under extreme low-data scenarios with the ultimate aim of reducing practical barriers to producing high spatial resolution land cover maps.
